\section{Two separately stable standard factors of unequal natural degrees}
\label{sec:unequal-standard-partners}

\paragraph{Status and shared provenance.}
HD-MS01--03, reserved question Q2eb, is a new derivation by
\texttt{/root/g34\_eightfolds}. Root supplied the exact target-one and
target-minus-one cycle formulas and the suggestion to retain both actual
natural quotient covers. The equal-degree Cauchy argument in HD-AST07
was a root derivation with an imprimitive contribution. The unequal
weighted-Cauchy loss, its all-target overlap argument and the source
consequences below are the present derivation. Complete author, root
and independent whole-proof, source-interface and code readings pass
without mathematical repair; their scopes are recorded separately. No previously accepted
equal-degree inequality is silently applied to unequal dimensions.

Let $n\ge m\ge4$, let $W_n,W_m$ be the ordinary deleted permutation
modules, and put $a=n-1$, $b=m-1$, $d=ab$. Write
$\nu(T)=\dim T-\max_\zeta\dim\ker(T-\zeta I)$.
For $x\in S_n,y\in S_m$, let $c_x,c_y$ count cycles including fixed
points and set $k=n-c_x$, $l=m-c_y$, $K=k+l$.

\subsection*{HD-MS01. A uniform all-target inequality}

\begin{lemma}[HD-MS01]
For every pair $(x,y)\ne(1,1)$,
\[
 \nu\bigl(W_n(x)\otimes W_m(y)\bigr)
 \ge F(k,l):=\frac{(m-3)k+(n-3)l+m}{2}.
 \tag{MS.1}
\]
The lemma holds for all $n\ge m\ge4$. There is no comparability
assumption on $n,m$, no restriction to target eigenvalue one, and
no coefficient-field hypothesis.
\end{lemma}

\begin{proof}
Put $c=c_x,d_0=c_y$, so $K=n+m-c-d_0$.
Use a common finite cyclic group for all eigenvalues. On the full
$n$-letter module, every cycle contributes at most one to each
eigenvalue multiplicity $p_\alpha$, so $p_\alpha\le c$ and
$\sum p_\alpha=n$. Removing one trivial line gives
\[
 \sum_\alpha u_\alpha^2
 =\sum_\alpha p_\alpha^2-2c+1\le(n-2)c+1,
 \qquad
 \sum_\beta v_\beta^2\le(m-2)d_0+1.
 \tag{MS.2}
\]
At any target $\zeta$, tensor multiplicity is
$\sum_\alpha u_\alpha v_{\zeta/\alpha}$.
Cauchy--Schwarz and the weighted arithmetic--geometric mean inequality
therefore give the all-target preliminary bound
\[
 \begin{split}
 \operatorname{mult}(\zeta)
 &\le\sqrt{((n-2)c+1)((m-2)d_0+1)}\\
 &\le\frac{(m-2)c+(n-2)d_0}{2}
       +\frac12\left(\frac{m-2}{n-2}+\frac{n-2}{m-2}\right),\\
 \nu&\ge B_0-E,\qquad
 B_0=\frac{(m-2)k+(n-2)l}{2},\quad
 E=\frac{(n-m)^2}{2(n-2)(m-2)}.
 \end{split}
 \tag{MS.3}
\]
Notice that $F=B_0-(K-m)/2$. The loss $E$ is retained rather
than discarded as in an equal-degree argument.

First dispose of a full-cycle factor. If $d_0=1$, then every
eigenvalue on $W_m(y)$ is simple. At any target the tensor
multiplicity is at most $a$, hence moving rank is at least
$a(b-1)$. If $c=1$, it is at most $b$ and moving rank is at
least $b(a-1)$. The maximum of $2F$ over
$0\le k\le a$, $0\le l\le b$ is
$(m-3)a+(n-3)b+m$. Subtracting it from twice these two moving
ranks gives respectively
\[
 2a(b-1)-\{(m-3)a+(n-3)b+m\}=m-2>0,
 \qquad
 2b(a-1)-\{(m-3)a+(n-3)b+m\}=2n-m-2>0.
 \tag{MS.4}
\]
Thus every target is handled when $\min(c,d_0)=1$.
For the rest of the proof assume $c,d_0\ge2$.

For a nonreal target, rationality of the standard tensor matrices
gives equal multiplicities at $\zeta$ and $\bar\zeta$, so moving
rank is at least $d/2$. If $K\le m$, the largest possible $2F$
occurs at $l=m-1,k=1$, and is $d-1$; this already proves the
claim. If $K\ge m+2E$, (MS.3) proves the claim.
In the only remaining interval put $T=K-m$, so $0<T<2E$.
Since $d_0\ge2$, we have $l\le m-2$, whence
\[
 2F=(m-3)K+(n-m)l+m
       \le n(m-2)+(m-3)T.
 \tag{MS.5}
\]
Writing $h=n-m\ge0$, the exact identity
\[
 \begin{split}
 &(h+1)(n-2)(m-2)-h^2(m-3)\\
 &\hspace{12mm}=h^2+(m-2)(m-1)h+(m-2)^2>0
 \end{split}
 \tag{MS.6}
\]
shows $(m-3)2E<n-m+1$. Thus (MS.5) is strictly less than
$n(m-2)+(n-m+1)=ab=d$. Rationality again gives the desired
bound. This proves every nonreal target without any condition
such as $(n-m)^2\le(n-2)(m-2)$.

It remains to handle the real targets $1,-1$.
Let $r_1,\ldots,r_c$ and $s_1,\ldots,s_{d_0}$ be the cycle
lengths. The permutation on the rectangular grid has
$C=\sum_{i,j}\gcd(r_i,s_j)$ cycles. Put
\[
 G=\frac{mc+nd_0}{2}-C\ge0.
 \tag{MS.7}
\]
The nonnegativity is the termwise inequality
$2\gcd(r_i,s_j)\le r_i+s_j$.
The identity $W_n\otimes W_m=P_n\otimes P_m-P_n-P_m+1$
gives target-one multiplicity $C-c-d_0+1$. Its moving rank is
exactly $B_0+G$. If $K\ge m$, this is at least $F$.
If $1\le K<m$, use
$\gcd(r,s)\le1+(r-1)(s-1)$ to get
$C\le cd_0+kl$. Consequently
\[
 G\ge\frac{mk+nl-4kl}{2}
   \ge\frac{mK-K^2}{2}
   =\frac{K(m-K)}2\ge\frac{m-K}{2}.
 \tag{MS.8}
\]
Here $n\ge m$ and $4kl\le K^2$ were used. This proves
the target-one assertion, including every small-index case.

For target minus one let $u,v$ count odd cycles in $x,y$.
The full-grid minus-one multiplicity and the corresponding
defect are
\[
 C_- =\sum_{\operatorname{lcm}(r_i,s_j)\ {\rm even}}
                \gcd(r_i,s_j),\qquad
 H=\frac{mc+nd_0}{2}-C_-.
 \tag{MS.9}
\]
The standard-pair minus-one multiplicity is
$C_--(c-u)-(d_0-v)$. Its moving rank is therefore
$B_0+H-u-v+1$. An odd--odd cycle pair contributes at least
one to $H$, an odd--even pair at least one half, and an
even--even pair a nonnegative amount. Hence
\[
 H\ge\frac{u d_0+v c}{2},\qquad
 H-u-v+1\ge1+\frac{u(d_0-2)+v(c-2)}2.
 \tag{MS.10}
\]
If $K\ge m$, the last expression is at least one because
$c,d_0\ge2$, and is more than the needed $(m-K)/2$.
If $K<m$, put $L=m-K=c+d_0-n>0$.
When $\min(c,d_0)=2$, one has $L\le2$, so the same lower
bound one is sufficient. When $c,d_0\ge3$, the length count
gives $u\ge2c-n$, $v\ge2d_0-m$, and thus
\[
 u+v\ge n-m+2L\ge2L,
 \qquad H-u-v+1\ge1+(u+v)/2\ge1+L>L/2.
 \tag{MS.11}
\]
These exhaust the real targets. An off-spectrum target has
full moving rank. Together with (MS.3)--(MS.6) they prove
(MS.1) for every target and every $n\ge m\ge4$.
\end{proof}

\subsection*{HD-MS02. Two actual quotient covers and every genuine mark}

\begin{theorem}[HD-MS02]
Suppose an actual finite coefficient image on
\[
 V=W_n\otimes W_m\otimes U\quad\hbox{up to scalar phases},
 \qquad q=\dim U\ge1,\qquad e=abq,
 \tag{MS.12}
\]
has two separately stable standard factors and actual transitive
natural quotient homomorphisms into $S_n,S_m$, with $n\ge m\ge4$.
Assume every nonscalar operator in the entire remaining partner
has normalized projective support at least $1/n$. For $q=1$
this condition is vacuous. If a connected sphere source has
$r\ge12$ genuine coefficient inertias and moving-rank sum $2e+6$,
then necessarily
\[
                       q(4m-2n+4)\le6.
 \tag{MS.13}
\]
In particular no such source exists when $n\le2m-2$.
The assertion allows every original coefficient field and scalar
phase. It requires no transitivity on the $nm$-letter product
action and no component-interchanging operator.
\end{theorem}

\begin{proof}
Put $C=(m-3)/2$, $D=(n-3)/2$, $B=m/2$.
At an inertia whose pair of natural projections is nonidentity,
HD-MS01 and spectral convolution with the partner give
\[
          \operatorname{rank}(1-T_i)
                  \ge q(Ck_i+Dl_i+B).
 \tag{MS.14}
\]
This holds for every scalar phase because HD-MS01 bounds every
eigenvalue target. If both projections are identity, a nonscalar
partner gives moving rank at least $e/n\ge qB$, since
\[
 2(n-1)(m-1)-mn=(n-2)(m-2)-2>0
 \quad(n\ge m\ge4).
 \tag{MS.15}
\]
If the partner is scalar as well, genuineness gives full actual
moving rank and hence the same lower bound. Thus (MS.14) is
valid for every original mark, including identity projections.
No twelve-mark assertion is made for either projected tuple.

The two genuine connected quotient covers separately give
$\sum k_i\ge2n-2$ and $\sum l_i\ge2m-2$.
Since $C,D,B>0$ and $r\ge12$, summing (MS.14) gives
\[
 \begin{split}
 \sum_i\operatorname{rank}(1-T_i)-2e
 &\ge q\{(m-3)(n-1)+(n-3)(m-1)+6m-2ab\}\\
 &=q(4m-2n+4).
 \end{split}
 \tag{MS.16}
\]
The exact budget proves (MS.13). If $n\le2m-2$, its right
side in (MS.16) is at least $8q>6$, a contradiction.

The nearby endpoints are kept distinct. At $n=2m-1$ or $n=2m$,
(MS.13) permits only $q=1$; at $n=2m+1$ it permits only
$q\le3$. These are necessary restrictions, not existence
statements. When $n\ge2m+2$ the displayed excess need not
be positive, and this theorem makes no universal exclusion there.
\end{proof}

\subsection*{HD-MS03. A second stable standard layer must be much smaller}

\begin{corollary}[HD-MS03]
Fix the complete recursive decomposition D-TD12 of an actual
finite primitive coefficient image in a rank-nine sixfold supplier.
Suppose a separately stable ordinary standard extraction layer
has natural degree $n\ge80$, maximal among all alternating
natural degrees in every layer of that same decomposition.
Then any other separately stable ordinary standard extraction
layer of natural degree $m\le n$ satisfies
\[
                              m\le\frac{n+1}{2}.
 \tag{MS.17}
\]
More precisely, if $m\ge8$ and $q$ is the dimension of the
entire complement of the two selected factors, (MS.13) holds.
The statement concerns separately stable whole layers, each
with one component; it does not select one factor from a
permuted lower-degree layer.
\end{corollary}

\begin{proof}
If $m\le7$, (MS.17) is immediate from $n\ge80$.
Otherwise both alternating natural degrees are at least eight.
The accepted standard-normalizer interface gives actual natural
homomorphisms into $S_n,S_m$, whose images contain $A_n,A_m$.
The scalar intersection of each standard image is trivial;
fixed sign-extension choices only change scalar phases.
Both natural actions are transitive. These homomorphisms are
pulled back from the genuine projective quotients of D-TD12;
no independent subgroup lifts of the later cores are assumed.

Both selected layer matrix algebras are stable. Reordering the
two tensor factors therefore leaves their entire complement,
including earlier and later extraction pieces, as an actual
finite projective partner. D-TD12 omits no residual space, and
HD-FT21 gives its support floor at least
$1/\max\{80,m_{*,U}-1\}\ge1/n$ by the same-decomposition
maximality. The empty partner is again harmless.
Here $e\ge(n-1)(m-1)>6$, so the accepted intermediate-cover
interface supplies the sphere, twelve genuine marks and budget
used by HD-MS02. If $m>(n+1)/2$, integrality gives
$n\le2m-2$, and that theorem is a contradiction.
The low-degree observation at the start does not add an
exceptional $A_6$ automorphism assertion to this corollary.
\end{proof}

\paragraph{Range, strictness, and remaining boundary.}
The local lemma HD-MS01 is universal in $n\ge m\ge4$.
An earlier sufficient Cauchy interval
$(n-m)^2\le(n-2)(m-2)$ is unnecessary because of (MS.6);
it is not the hypothesis of the frozen lemma. The new universal
source-exclusion range is exactly the sufficient range
$n\le2m-2$ proved above. The neighboring bounds in (MS.13)
are not reversed into existence or exclusion claims.

For example the parameter configurations $n=2M,m=M+1$,
$M\ge40$, have two different separately stable standard layers
and source excess at least $8q$. They are outside the same-orbit
permuted-standard hypothesis of HD-AST14--15 and were retained
by the single-standard alternative of HD-ML01. Thus this is a
new excluded parameter class, not a numerical restatement of
that same-orbit theorem. No actual supplier with those parameters
is asserted to exist. Smaller standard partners, permuted
lower-degree layers, nonstandard partners and any omitted
partner factor remain outside a broader conclusion unless their
own exact interfaces are proved. There is no new field descent,
polarization, inducing-order bound or detecting cycle construction.
